% Problem record op_ac5fa4581b04f4c2. This TeX file is derived from
% database/problems_json/op_ac5fa4581b04f4c2.json by scripts/sync-tex.mjs; edit the JSON record
% and rerun the script rather than editing this file.
\section{Additivity of the entanglement of purification}
\label{sec:op_ac5fa4581b04f4c2}

\paragraph{Problem.}
Is the entanglement of purification additive on tensor products?  For a
bipartite density operator $\rho_{AB}$ on finite-dimensional Hilbert spaces
$A$ and $B$, define its entanglement of purification by
\begin{equation}
  E_P(A{:}B)_\rho
  :=\min_{\psi}\ S(AA')_\psi,
  \label{eq:epur-definition}
\end{equation}
where the minimum ranges over all pure states $\psi$ on $AA'BB'$ with
purifying systems $A'$, $B'$ of arbitrary finite dimension such that
$\operatorname{Tr}_{A'B'}\psi=\rho_{AB}$, and $S$ denotes the von Neumann
entropy.  Taking the product of optimal purifications gives
$E_P(\rho\otimes\sigma)\leq E_P(\rho)+E_P(\sigma)$, so by Fekete's lemma for
subadditive sequences the normalized values $\frac1nE_P(\rho^{\otimes n})$
converge to their infimum rather than decrease monotonically, and the
regularization
\begin{equation}
  E_P^\infty(\rho)
  :=\lim_{n\to\infty}\frac1nE_P\bigl(\rho^{\otimes n}\bigr)
  =\inf_{n\geq1}\frac1nE_P\bigl(\rho^{\otimes n}\bigr)
  \label{eq:epur-regularization}
\end{equation}
is well defined.  The question is whether regularization is unnecessary:
does
\begin{equation}
  E_P(\rho\otimes\sigma)=E_P(\rho)+E_P(\sigma)
  \quad\text{for all finite-dimensional bipartite }\rho\text{ and }\sigma,
  \label{eq:epur-additivity}
\end{equation}
hold?  Applied inductively to the pairs $(\rho,\rho^{\otimes(n-1)})$,
Eq.~\eqref{eq:epur-additivity} would give $E_P(\rho^{\otimes n})=nE_P(\rho)$
and hence make the regularization of Eq.~\eqref{eq:epur-regularization}
collapse to $E_P^\infty(\rho)=E_P(\rho)$ for every $\rho$; whether such collapse on
single-state tensor powers would already force Eq.~\eqref{eq:epur-additivity}
for arbitrary pairs is not known.  Since Eq.~\eqref{eq:epur-additivity} can
only fail through strict subadditivity for some pair, the simplest
potential witness is
\begin{equation}
  E_P\bigl(\rho^{\otimes2}\bigr)<2E_P(\rho),
  \label{eq:epur-strict-subadditivity}
\end{equation}
a violating pair with $\sigma=\rho$; no reduction from a violation with
distinct $\rho\neq\sigma$ to such a same-state violation is known.
The sharpest openly posed subcase concerns two-qubit classical states
$\rho=\sum_{i,j\in\{0,1\}}p_{ij}|ij\rangle\langle ij|$ at the von Neumann
order: for this family the R\'enyi generalizations of
Eq.~\eqref{eq:epur-definition} are settled non-additive for every order
$\alpha\in[0,1)$ and additive for every $\alpha\in[2,\infty]$, while the
interval $\alpha\in[1,2)$, which includes the von Neumann case $\alpha=1$
of Eqs.~\eqref{eq:epur-additivity}--\eqref{eq:epur-strict-subadditivity},
remains open.

\subsection*{Status}

Solved

\subsection*{Source}

The additivity question accompanies the introduction of the measure: Terhal,
Horodecki, Leung, and DiVincenzo define $E_P$ as in
Eq.~\eqref{eq:epur-definition} and prove that the regularization of
Eq.~\eqref{eq:epur-regularization} equals the entanglement cost of creating
the state from maximally entangled states with negligible communication
\sourcecite{ref:epur-terhal}{THLD02}, making its possible collapse to $E_P$
the natural open question, and Chen and Winter investigate the failure
alternative of Eq.~\eqref{eq:epur-strict-subadditivity} explicitly
\sourcecite{ref:epur-chen-winter}{CW12}.

\subsection*{Progress}

\begin{itemize}
  \item Terhal, Horodecki, Leung, and DiVincenzo introduced $E_P$ in the form of
Eq.~\eqref{eq:epur-definition}, proved that $E_P^\infty$ equals the
entanglement cost of creating the state from maximally entangled states with
negligible communication, and showed that $\frac12I(A{:}B)_\rho$ and the
regularized classical mutual information --- maximized over local
measurements, $I_c^\infty$ in THLD02's notation --- lower-bound
$E_P^\infty(\rho)$; they also reported numerics on Werner states
\sourcecite{ref:epur-terhal}{THLD02}.

  \item Chen and Winter proved that $E_P^\infty(\rho)-S(\rho)$ is convex, via a new
asymptotic preparation protocol, and gave numerical evidence for the failure of
Eq.~\eqref{eq:epur-additivity} on two-qubit Werner states
\sourcecite{ref:epur-chen-winter}{CW12}.

  \item Bagchi and Pati confirmed that $E_P$ is subadditive on tensor products, so
that Eq.~\eqref{eq:epur-additivity} can only fail in the direction of
Eq.~\eqref{eq:epur-strict-subadditivity}, identified classes of states on
which $E_P$ is additive on tensor products, and proved polygamy of $E_P$ for
tripartite pure states \sourcecite{ref:epur-bagchi-pati}{BP15}.

  \item Baghali Khanian proved that the optimal visible-compression rate of mixed
states is expressed through the entanglement of purification, with
exponential fidelity decay below a regularized R\'enyi threshold, proved
additivity in Eq.~\eqref{eq:epur-additivity} only for a variation of
extendible states, and recorded that the additivity of $E_P$ remains unknown
\sourcecite{ref:epur-khanian}{KB25}.

  \item Faraji and Baghali Khanian proved additivity of the R\'enyi-2 entanglement of
purification for the states associated with transpose-depolarizing and
depolarizing channels and their complements, via multiplicativity of the
constrained output 2-norm; their discussion records that for Werner states
Eq.~\eqref{eq:epur-additivity} is known at the parameter values
$t\in\{0,\,1/(1-d),\,1/(1+d)\}$, attributed to Terhal et al. and to
Christandl and Winter, and that the putative $d=2$ violation rests on
numerics, leaving the Werner R\'enyi additivity question open
\sourcecite{ref:epur-faraji}{FBK26},
\sourcecite{ref:epur-christandl-winter}{CW05},
\sourcecite{ref:epur-terhal}{THLD02}.

  \item Negari and Baghali Khanian solved the one-copy optimization for two-qubit
classical states at every R\'enyi order, proved non-additivity for every
$\alpha\in[0,1)$, witnessed by two copies of a single state, and additivity
for every $\alpha\in[2,\infty]$; they state that the interval
$\alpha\in[1,2)$, including the von Neumann case $\alpha=1$ of
Eq.~\eqref{eq:epur-additivity}, remains open, with additivity conjectured
there \sourcecite{ref:epur-negari}{NBK26}.

  \item Baghali Khanian and Negari disprove Eq.~\eqref{eq:epur-additivity} using the two-qubit Werner state $W(f)$ with singlet weight $f$ and equal weight $(1-f)/3$ on each triplet. Proposition~IV.3 identifies its one-copy optimum with the entanglement $C(f)$ of the canonical purification across $AA':BB'$ for $1/64\leq f\leq 1/4$. Proposition~III.1 gives $E_P^\infty(W(f))<C(f)$ for $0<f<f_*$, where $f_*\approx 0.01708596$ is the root of Eq.~(B.3). At $f=1/64$, Eq.~(9) of the paper applies the Chen--Winter ensemble bound to $W(f)=p W(0)+(1-p)W(t)$ with $p=9/256$ and $t=f/(1-p)=4/247$. The resulting upper bound $V_f(p)$ gives, in bits,
\begin{equation}
  E_P^\infty(W(f))\leq V_f(p)<0.848959<0.849252<E_P(W(f)).
  \label{eq:epur-werner-gap}
\end{equation}
The positive gap persists for every $1/64\leq f<f_*$ \sourcecite{ref:epur-werner-nonadditivity}{BKN26}, \sourcecite{ref:epur-chen-winter}{CW12}
(\href{https://github.com/Naixu-Guo/quantum-open-problems/issues/132}{GitHub \#132}).
\end{itemize}

\subsection*{References}

\begin{enumerate}[label={},leftmargin=4.8em,labelsep=0.5em]
  \footnotesize
  \item[\textup{[THLD02]}]\label{ref:epur-terhal}
  B. M. Terhal, M. Horodecki, D. W. Leung, and D. P. DiVincenzo,
  ``The Entanglement of Purification,'' \emph{Journal of Mathematical Physics}
  \textbf{43}(9), 4286--4298 (2002).
  \href{https://arxiv.org/abs/quant-ph/0202044}{arXiv:quant-ph/0202044}.

  \item[\textup{[CW12]}]\label{ref:epur-chen-winter}
  J. Chen and A. Winter,
  ``Non-Additivity of the Entanglement of Purification (Beyond Reasonable
  Doubt),'' arXiv preprint (2012).
  \href{https://arxiv.org/abs/1206.1307}{arXiv:1206.1307}.

  \item[\textup{[BP15]}]\label{ref:epur-bagchi-pati}
  S. Bagchi and A. K. Pati,
  ``Monogamy, Polygamy, and Other Properties of Entanglement of
  Purification,'' \emph{Physical Review A} \textbf{91}, 042323 (2015).
  \href{https://arxiv.org/abs/1502.01272}{arXiv:1502.01272}.

  \item[\textup{[KB25]}]\label{ref:epur-khanian}
  Z. Baghali Khanian,
  ``Strong Converse Bounds for Compression of Mixed States,'' arXiv preprint
  (2025), version~2, revised 21 April 2025.
  \href{https://arxiv.org/abs/2206.09415v2}{arXiv:2206.09415v2}.

  \item[\textup{[FBK26]}]\label{ref:epur-faraji}
  S. Faraji and Z. Baghali Khanian,
  ``Additivity Results for the R\'enyi-2 Entanglement of Purification,''
  arXiv preprint (2026).
  \href{https://arxiv.org/abs/2605.15439}{arXiv:2605.15439}.

  \item[\textup{[NBK26]}]\label{ref:epur-negari}
  A.-R. Negari and Z. Baghali Khanian,
  ``R\'enyi Entanglement of Purification Is Non-additive,'' arXiv preprint
  (2026).
  \href{https://arxiv.org/abs/2608.28897}{arXiv:2608.28897}.

  \item[\textup{[CW05]}]\label{ref:epur-christandl-winter}
  M. Christandl and A. Winter,
  ``Uncertainty, Monogamy, and Locking of Quantum Correlations,''
  \emph{IEEE Transactions on Information Theory} \textbf{51}(9), 3159--3165
  (2005).
  \href{https://arxiv.org/abs/quant-ph/0501090}{arXiv:quant-ph/0501090}.

  \item[\textup{[BKN26]}]\label{ref:epur-werner-nonadditivity}
  Z. Baghali Khanian and A.-R. Negari,
  ``Entanglement of Purification for Werner States: Canonical Purification and Nonadditivity,'' arXiv preprint, version~1 (2026).
  \href{https://arxiv.org/abs/2610.03630v1}{arXiv:2610.03630v1}.
\end{enumerate}

\subsection*{Comment}

The regularized gap proves strict subadditivity for a finite block. Set $W=W(1/64)$ and $a_n=E_P(W^{\otimes n})$. If $a_{2^{j+1}}=2a_{2^j}$ for every $j\geq 0$, then $a_{2^j}/2^j=a_1$ for every $j$. Equation~\eqref{eq:epur-regularization} makes these ratios converge to $E_P^\infty(W)<a_1$. Hence Eq.~\eqref{eq:epur-strict-subadditivity} holds for some finite-dimensional $\rho=W^{\otimes 2^j}$. Whether $j=0$ suffices, giving a violation on two copies of the original Werner state, remains open. The Werner witnesses are separable and have correlations that cannot be diagonalized in a local product basis. The von Neumann question for classical two-qubit states also remains open \sourcecite{ref:epur-werner-nonadditivity}{BKN26}, Sec.~V, \sourcecite{ref:epur-negari}{NBK26}.

\subsection*{Field}

Quantum Resource Theory

\subsection*{Topic}

Entanglement measures; Additivity and regularization

\subsection*{ID}

\texttt{op\_ac5fa4581b04f4c2}
